\section{Limitation}
% 尽管 TORQ 展现出优异的精度与效率权衡，仍存在若干值得进一步探索的局限性：首先，块内旋转当前聚焦于一阶统计（log 幅值期望）的匹配，未能充分利用激活分布的高阶统计信息（如峰度、偏度），在部分激活分布极度非对称或多峰的场景下，可能无法完全释放 FP4 的表达潜力；其次，本文采用固定块大小（如 32）进行实验，而不同模型层、不同任务的激活特性存在显著差异，自适应块大小的动态选择机制及其对量化精度与硬件效率的影响尚未深入研究；最后，虽然我们在 MoE 模型上验证了方法的有效性，但未针对 Expert 间激活分布的异质性设计专属优化策略，对于超大规模 MoE 模型（如 Expert 数量超千级），如何进一步提升跨 Expert 量化一致性仍需探索。这些局限也为未来工作指明了方向，例如结合高阶统计的自适应旋转设计、块大小与模型特性的联合优化等。
Although TORQ demonstrates an excellent trade-off between accuracy and efficiency, there are several limitations that warrant further exploration: First, the current intra-block rotation focuses on matching first-order statistics (the expectation of log magnitudes) and fails to fully utilize the higher-order statistical information of the activation distribution (such as kurtosis and skewness). In scenarios where some activation distributions are extremely asymmetric or multimodal, it may not fully unleash the expressive potential of FP4. Second, this paper uses a fixed block size (e.g., 32) for experiments, while the activation characteristics of different model layers and different tasks vary significantly. The dynamic selection mechanism for adaptive block sizes and its impact on quantization accuracy and hardware efficiency have not been thoroughly studied. Finally, although we verified the effectiveness of the method on MoE models, we did not design specific optimization strategies for the heterogeneity of activation distributions among Experts. For ultra-large-scale MoE models (e.g., with more than a thousand Experts), how to further improve cross-Expert quantization consistency remains to be explored. These limitations also point out directions for future work, such as the design of adaptive rotation combined with higher-order statistics and the joint optimization of block size and model characteristics.
